The proved capacity drift law gives $w=-1+\tfrac13 d\log N/d\log a$.
With an area capacity $N\propto R^2$, a comoving radius gives $w=-1/3$.
A Hubble radius with separately conserved dust instead gives $w=0$ and
$H^2/H_0^2=a^{-3}$. Retaining mechanical pressure $-\rho$ on that branch
requires matter exchange and a constant dark fraction, with no transition
from matter domination to acceleration.

For the declared future-event-horizon member with coefficient one,
$x=\sqrt{\Omega_{\mathrm{dark}}}$ obeys the exact algebraic equation
\[
 \frac{x^6(1+x)^3}{(1-x)(1+2x)^8}
 =a^3\frac{x_0^6(1+x_0)^3}{(1-x_0)(1+2x_0)^8}.
\]
Monotonicity and rational bisection replace an ordinary differential equation
solver. The future limit has positive finite $H$, and the event-horizon
boundary condition follows. With the supplied $H_0=67.4\,\mathrm{km\,s^{-1}
Mpc^{-1}}$ and $\Omega_m=0.315$, all fifteen correlated primary chronometer
rows give $\chi^2\simeq6.57717$. The same comparison gives $6.19134$ for
fixed capacity, $21.54903$ for comoving area and $70.48705$ for separately
conserved Hubble area. The alternative stellar-population reductions use
the same observations. These are conditional capacity-history comparisons;
neither the cosmic amplitude nor a physical register-to-horizon map is derived.
